\documentclass[11pt,a4paper]{article}

\usepackage[margin=1in]{geometry}
\usepackage{amsmath}
\usepackage{booktabs}
\usepackage{array}
\usepackage{tabularx}
\usepackage{hyperref}
\usepackage{enumitem}

\hypersetup{
  colorlinks=true,
  linkcolor=black,
  urlcolor=blue
}

\setlist[itemize]{leftmargin=*,nosep}
\renewcommand{\arraystretch}{1.2}
\newcolumntype{Y}{>{\raggedright\arraybackslash}X}

\begin{document}

\title{Bucky++ Protection Test Plan}
\date{}
\maketitle
\tableofcontents

\section*{Scope}
\addcontentsline{toc}{section}{Scope}

Passive ranges below use the KiCad-installed passive tolerances visible in the project: resistors $\pm 1\%$ and timing capacitors $\pm 10\%$. They are intended as board-level expectations before bench test. Internal IC accuracy is not added unless it is already embedded in the review math.

\section{TPS259830ONRGER}

\subsection{Theoretical Values}

\texttt{IC7}, \texttt{IC9}, \texttt{IC8}, and \texttt{IC10} implement rail e-fuse protection on \texttt{5V}, \texttt{12V}, \texttt{15V}, and \texttt{24V}.

For each rail the review math is:

\begin{align}
V_{UV} &= 1.2 \cdot \frac{R_1 + R_2 + R_3}{R_2 + R_3} \\
V_{OV} &= 1.2 \cdot \frac{R_1 + R_2 + R_3}{R_3} \\
I_{LIMIT} &= \frac{1460}{R_{ILIM}} + 0.11 \\
t_{ITIMER} &= 470 \cdot C_{ITIMER} \\
t_{RETRY} &\approx \frac{\left(C_{RETRY\_DLY} \cdot 1000 + 4\right) \cdot 46.83}{1000}\,\mathrm{ms}
\end{align}

Installed values from the KiCad project:

\begin{itemize}
\item \textbf{5 V rail, IC7:} $R11 = 1\,\mathrm{M\Omega}$, $R12 = 142\,\mathrm{k\Omega}$, $R13 = 285\,\mathrm{k\Omega}$, $R10 = 150\,\mathrm{\Omega}$, $C15 = 22\,\mathrm{nF}$, $C22 = 47\,\mathrm{nF}$
\item \textbf{12 V rail, IC9:} $R21 = 1\,\mathrm{M\Omega}$, $R22 = 39\,\mathrm{k\Omega}$, $R23 = 97\,\mathrm{k\Omega}$, $R25 = 150\,\mathrm{\Omega}$, $C18 = 22\,\mathrm{nF}$, $C19 = 47\,\mathrm{nF}$
\item \textbf{15 V rail, IC8:} $R15 = 1\,\mathrm{M\Omega}$, $R16 = 24\,\mathrm{k\Omega}$, $R17 = 78\,\mathrm{k\Omega}$, $R19 = 150\,\mathrm{\Omega}$, $C11 = 22\,\mathrm{nF}$, $C31 = 47\,\mathrm{nF}$
\item \textbf{24 V rail, IC10:} $R27 = 1\,\mathrm{M\Omega}$, $R28 = 9\,\mathrm{k\Omega}$, $R29 = 49\,\mathrm{k\Omega}$, $R31 = 150\,\mathrm{\Omega}$, $C10 = 22\,\mathrm{nF}$, $C14 = 47\,\mathrm{nF}$
\end{itemize}

Calculated thresholds:

\paragraph{5 V rail (IC7)}
\begin{align*}
V_{UV} &= 1.2 \cdot \frac{1000 + 142 + 285}{142 + 285} = 4.01\,\mathrm{V} \\
V_{OV} &= 1.2 \cdot \frac{1427}{285} = 6.01\,\mathrm{V}
\end{align*}
Tolerance range: $V_{UV} = 3.95\,\mathrm{V}$ to $4.07\,\mathrm{V}$, $V_{OV} = 5.91\,\mathrm{V}$ to $6.11\,\mathrm{V}$.

\paragraph{12 V rail (IC9)}
\begin{align*}
V_{UV} &= 1.2 \cdot \frac{1000 + 39 + 97}{39 + 97} = 10.02\,\mathrm{V} \\
V_{OV} &= 1.2 \cdot \frac{1136}{97} = 14.05\,\mathrm{V}
\end{align*}
Tolerance range: $V_{UV} = 9.85\,\mathrm{V}$ to $10.20\,\mathrm{V}$, $V_{OV} = 13.80\,\mathrm{V}$ to $14.31\,\mathrm{V}$.

\paragraph{15 V rail (IC8)}
\begin{align*}
V_{UV} &= 1.2 \cdot \frac{1000 + 24 + 78}{24 + 78} = 12.96\,\mathrm{V} \\
V_{OV} &= 1.2 \cdot \frac{1102}{78} = 16.95\,\mathrm{V}
\end{align*}
Tolerance range: $V_{UV} = 12.73\,\mathrm{V}$ to $13.20\,\mathrm{V}$, $V_{OV} = 16.64\,\mathrm{V}$ to $17.27\,\mathrm{V}$.

\paragraph{24 V rail (IC10)}
\begin{align*}
V_{UV} &= 1.2 \cdot \frac{1000 + 9 + 49}{9 + 49} = 21.89\,\mathrm{V} \\
V_{OV} &= 1.2 \cdot \frac{1058}{49} = 25.91\,\mathrm{V}
\end{align*}
Tolerance range: $V_{UV} = 21.48\,\mathrm{V}$ to $22.31\,\mathrm{V}$, $V_{OV} = 25.42\,\mathrm{V}$ to $26.41\,\mathrm{V}$.

\paragraph{Common settings}
\begin{align*}
I_{LIMIT} &= \frac{1460}{150} + 0.11 = 9.84\,\mathrm{A}
\end{align*}
Tolerance range: $9.75\,\mathrm{A}$ to $9.94\,\mathrm{A}$.

\begin{itemize}
\item $ITIMER$ is intentionally left open, so the overcurrent blanking interval is the minimum setting.
\item $C_{RETRY\_DLY} = 22\,\mathrm{nF}$ gives $t_{RETRY} \approx 1030\,\mathrm{ms}$, with a passive-only range of about $927\,\mathrm{ms}$ to $1133\,\mathrm{ms}$.
\item $C_{NRETRY} = 47\,\mathrm{nF}$ is set for $16$ retries per the TI table used in the review flow.
\end{itemize}

\subsection{Testing Plan}

\begin{itemize}
\item Use a programmable bench supply on the rail input and an electronic load on the protected output.
\item Sweep the input voltage down until the rail disconnects to measure the UV trip point.
\item Sweep the input voltage up until the rail disconnects to measure the OV trip point.
\item At nominal input voltage, increase load current until current limiting starts; record the steady current limit.
\item Apply a hard short or low-ohm overload and use an oscilloscope on \texttt{FAULT}, \texttt{PG}, and output voltage to measure retry delay and retry count.
\item Confirm no nuisance trip during normal startup at the intended load.
\end{itemize}

\subsection{Results Tables}

\subsubsection{Thresholds}
\small
\begin{tabularx}{\textwidth}{p{1.2cm}p{2.2cm}p{2.2cm}p{2.2cm}p{1.8cm}p{1.8cm}p{1.8cm}Y}
\toprule
Rail & UV theory & OV theory & OC theory & Measured UV & Measured OV & Measured OC & Notes \\
\midrule
5V  & $4.01\,(3.95\ldots4.07)\,\mathrm{V}$ & $6.01\,(5.91\ldots6.11)\,\mathrm{V}$ & $9.84\,(9.75\ldots9.94)\,\mathrm{A}$ & & & & \\
12V & $10.02\,(9.85\ldots10.20)\,\mathrm{V}$ & $14.05\,(13.80\ldots14.31)\,\mathrm{V}$ & $9.84\,(9.75\ldots9.94)\,\mathrm{A}$ & & & & \\
15V & $12.96\,(12.73\ldots13.20)\,\mathrm{V}$ & $16.95\,(16.64\ldots17.27)\,\mathrm{V}$ & $9.84\,(9.75\ldots9.94)\,\mathrm{A}$ & & & & \\
24V & $21.89\,(21.48\ldots22.31)\,\mathrm{V}$ & $25.91\,(25.42\ldots26.41)\,\mathrm{V}$ & $9.84\,(9.75\ldots9.94)\,\mathrm{A}$ & & & & \\
\bottomrule
\end{tabularx}
\normalsize

\vspace{0.8em}

\small
\begin{tabularx}{\textwidth}{p{1.2cm}p{2.6cm}p{2.2cm}p{2.2cm}p{2.0cm}Y}
\toprule
Rail & Retry delay theory & Retry count theory & Measured delay & Measured count & Notes \\
\midrule
5V  & $1.03\,(0.93\ldots1.13)\,\mathrm{s}$ & $16$ & & & \\
12V & $1.03\,(0.93\ldots1.13)\,\mathrm{s}$ & $16$ & & & \\
15V & $1.03\,(0.93\ldots1.13)\,\mathrm{s}$ & $16$ & & & \\
24V & $1.03\,(0.93\ldots1.13)\,\mathrm{s}$ & $16$ & & & \\
\bottomrule
\end{tabularx}
\normalsize

\section{ADM1270ACPZ}

\subsection{Theoretical Values}

\texttt{U4} protects the battery input path against UV, OV, OC, and SOA stress through foldback.

The review math is:

\begin{align}
V_{ISET} &= 3.6 \cdot \frac{R_{BOTTOM}}{R_{TOP} + R_{BOTTOM}} \\
V_{SENSE} &= \frac{V_{ISET}}{40} \\
I_{LIMIT} &= \frac{V_{SENSE}}{R_{SHUNT}} \\
V_{UV} &= 1.0 \cdot \frac{R_1 + R_2 + R_3}{R_2 + R_3} \\
V_{OV} &= 1.0 \cdot \frac{R_1 + R_2 + R_3}{R_3} \\
V_{OUT,FLB} &= V_{ISET} \cdot \frac{R_A + R_B}{R_B}
\end{align}

Installed values from the project and the review note:

\begin{itemize}
\item $R39 = 1\,\mathrm{m\Omega} \pm 1\%$ shunt
\item $V_{ISET}$ divider $= 3.6\,\mathrm{k\Omega} / 2.3\,\mathrm{k\Omega}$
\item UV/OV divider $= 1\,\mathrm{M\Omega} / 15\,\mathrm{k\Omega} / 37\,\mathrm{k\Omega}$
\item Foldback divider $= 1\,\mathrm{M\Omega} / 133\,\mathrm{k\Omega}$
\item $C23 = 100\,\mathrm{nF}$ and $C32 = 100\,\mathrm{nF}$ for timer and retry-off timing
\end{itemize}

Calculated thresholds:

\begin{align*}
V_{ISET} &= 3.6 \cdot \frac{2.3}{3.6 + 2.3} = 1.403\,\mathrm{V} \\
V_{SENSE} &= \frac{1.403}{40} = 35.08\,\mathrm{mV} \\
I_{LIMIT} &= \frac{35.08\,\mathrm{mV}}{1\,\mathrm{m\Omega}} = 35.08\,\mathrm{A} \\
V_{UV} &= 1.0 \cdot \frac{1000 + 15 + 37}{15 + 37} = 20.23\,\mathrm{V} \\
V_{OV} &= 1.0 \cdot \frac{1000 + 15 + 37}{37} = 28.43\,\mathrm{V} \\
V_{OUT,FLB} &= 1.403 \cdot \frac{1000 + 133}{133} = 11.96\,\mathrm{V}
\end{align*}

Passive-only ranges:

\begin{itemize}
\item $V_{ISET} = 1.386\,\mathrm{V}$ to $1.421\,\mathrm{V}$
\item $V_{SENSE} = 34.66\,\mathrm{mV}$ to $35.51\,\mathrm{mV}$
\item $I_{LIMIT} = 34.31\,\mathrm{A}$ to $35.87\,\mathrm{A}$
\item $V_{UV} = 19.85\,\mathrm{V}$ to $20.62\,\mathrm{V}$
\item $V_{OV} = 27.89\,\mathrm{V}$ to $28.99\,\mathrm{V}$
\item $V_{OUT,FLB} = 11.60\,\mathrm{V}$ to $12.32\,\mathrm{V}$
\item $C_{TIMER} = 100\,\mathrm{nF}$ gives nominal fault timing of about $10\,\mathrm{ms}$
\item $C_{TIMER\_OFF} = 100\,\mathrm{nF}$ gives nominal retry-off time of about $200\,\mathrm{ms}$
\item With $\pm 10\%$ timing capacitors, timing-only expectation is about $9\,\mathrm{ms}$ to $11\,\mathrm{ms}$ for fault timing and $180\,\mathrm{ms}$ to $220\,\mathrm{ms}$ for retry-off time
\end{itemize}

\subsection{Testing Plan}

\begin{itemize}
\item Drive the battery input with a programmable supply and place an electronic load on \texttt{VBatt\_safe}.
\item Sweep the input voltage down and up slowly to measure the UV and OV trip points.
\item At nominal battery voltage, increase the load current until the current limiter engages; record the limited current.
\item Create a controlled overload or short so that $V_{OUT}$ is forced downward and measure the point where foldback starts reducing the current limit; target is near $12\,\mathrm{V}$.
\item Use an oscilloscope on output voltage, gate drive, and fault-related pins to measure fault-on time and retry-off time during a hard overload.
\end{itemize}

\subsection{Results Table}

\small
\begin{tabularx}{\textwidth}{p{2.6cm}p{4.0cm}p{2.4cm}p{2.0cm}Y}
\toprule
Feature & Theory & Measured & Pass/Fail & Notes \\
\midrule
UV trip & $20.23\,(19.85\ldots20.62)\,\mathrm{V}$ & & & \\
OV trip & $28.43\,(27.89\ldots28.99)\,\mathrm{V}$ & & & \\
Current limit & $35.08\,(34.31\ldots35.87)\,\mathrm{A}$ & & & \\
Foldback start & $11.96\,(11.60\ldots12.32)\,\mathrm{V}$ & & & \\
Fault timing & $10\,\mathrm{ms}$ nominal, $9\ldots11\,\mathrm{ms}$ from $C$ tolerance & & & \\
Retry-off timing & $200\,\mathrm{ms}$ nominal, $180\ldots220\,\mathrm{ms}$ from $C$ tolerance & & & \\
\bottomrule
\end{tabularx}
\normalsize

\section{LTC4372CMS8\#TRPBF}

\subsection{Theoretical Values}

\texttt{IC701}, \texttt{IC702}, \texttt{IC4}, and \texttt{IC5} implement reverse-current protection on \texttt{5V}, \texttt{12V}, \texttt{15V}, and \texttt{24V}.

The review note for the selected MOSFET states:

\begin{itemize}
\item $V_{GS(th)} = 2.33\,\mathrm{V}$ to $3.15\,\mathrm{V}$
\item $R_{DS(on)} = 1.30\,\mathrm{m\Omega}$ to $1.63\,\mathrm{m\Omega}$
\end{itemize}

Each protected path uses a back-to-back MOSFET pair, so:

\begin{align}
R_{PATH} &\approx 2 \cdot R_{DS(on)} = 2.60\,\mathrm{m\Omega} \text{ to } 3.26\,\mathrm{m\Omega} \\
V_{DROP} &= I_{TEST} \cdot R_{PATH}
\end{align}

Examples:

\begin{itemize}
\item At $10\,\mathrm{A}$, expected forward drop is $26\,\mathrm{mV}$ to $32.6\,\mathrm{mV}$.
\item At $20\,\mathrm{A}$, expected forward drop is $52\,\mathrm{mV}$ to $65.2\,\mathrm{mV}$.
\item Reverse-current expectation is functional, not adjustable: no sustained back-feed from output to input when the output is driven externally above the input or when the input source is removed.
\item \texttt{SHDN} is expected to open the path on command.
\end{itemize}

\subsection{Testing Plan}

\begin{itemize}
\item Power each rail from its normal input and load the protected output to a chosen current point; measure voltage drop across the protection stage.
\item Turn the input supply off, then drive the output connector from a bench supply with current limit enabled; verify that the input side does not get back-powered.
\item Toggle \texttt{SHDN\_CTRL} and confirm that the path opens cleanly and remains reverse-blocking.
\item If surge equipment is available, run a controlled non-destructive clamp test on the TVS network as a separate optional test.
\end{itemize}

\subsection{Results Table}

\small
\begin{tabularx}{\textwidth}{p{1.2cm}p{3.0cm}p{2.8cm}p{2.4cm}p{2.2cm}p{2.2cm}p{1.8cm}Y}
\toprule
Rail & Forward-drop theory & Reverse-current theory & SHDN theory & Measured forward drop & Measured reverse current & Pass/Fail & Notes \\
\midrule
5V  & $V_{DROP} = I_{TEST} \cdot (2.60\ldots3.26)\,\mathrm{m\Omega}$ & No sustained back-feed & Path opens on command & & & & \\
12V & $V_{DROP} = I_{TEST} \cdot (2.60\ldots3.26)\,\mathrm{m\Omega}$ & No sustained back-feed & Path opens on command & & & & \\
15V & $V_{DROP} = I_{TEST} \cdot (2.60\ldots3.26)\,\mathrm{m\Omega}$ & No sustained back-feed & Path opens on command & & & & \\
24V & $V_{DROP} = I_{TEST} \cdot (2.60\ldots3.26)\,\mathrm{m\Omega}$ & No sustained back-feed & Path opens on command & & & & \\
\bottomrule
\end{tabularx}
\normalsize

\end{document}
